\section{Arithmetic certificates and the iPRA proof}
\label{sec:arith}

\subsection{Numerical coding}

Here is one sufficient coding; the result does not depend on the
particular pairing convention. Encode a nonnegative integer \(z\)
by \(2z\) and a negative integer \(z\) by \(-2z-1\).
Encode a rational by a pair consisting of a signed numerator and
a positive denominator. A denominator can be stored as its
predecessor, so every such code has a valid positive denominator.
Reduction to lowest terms is unnecessary. Rational equality and
order are signed cross-multiplication tests.

For pairs, one may use
\(\langle u,v\rangle=2^u(2v+1)-1\). Its projections are bounded
searches on the input code. Fixed-length tuples and finite lists
are iterated pairs with a length tag; list lookup and finite
folds are primitive recursive. Matrices are lists of rows with
fixed dimensions. Invalid list lengths can be assigned a default
value outside guarded assertions. All evaluators are total
numerical functions.

Signed addition, multiplication, absolute value, finite maxima,
and rational division with a positive or certified nonzero
denominator are primitive recursive. The dyadic floor in
\eqref{eq:algorithm} is integer division with the correct
negative-numerator case. The box projection is a finite sequence
of rational order tests. These are decisions about rational
codes, not decisions about equality or order of arbitrary reals.

If \(u,v\) are regular names, equality can be expressed by
\[
 \operatorname{Eq}(u,v)\ :\!\iff\
 \forall p\ \norm{u(p)-v(p)}\le2^{1-p}.
\]
This expresses equality of their represented values. Refining
the observation index using \eqref{eq:regular} proves symmetry,
transitivity, and compatibility with the operations below.
A norm bound against a rational \(c\ge0\) is expressed by
\(\forall p\,\norm{u(p)}\le c+2^{-p}\).
Sum and difference names are obtained by shifting each input
index by one. A product name uses rational bounds on its
factors and a sufficiently large fixed precision shift.
Every product in the proof has such a bound: box radii, \(M\),
\(K\), \(2\), or the finite sum bound \(4\).

\subsection{Raw input certificates}\label{sec:raw}

For explicit first-order input, fix the primitive recursive
functions
\[
 \Hhat(q,t),\quad \Jhat(q,t),\quad
 \mu_0(p),\quad \nu_0(p).
\]
Their arguments are rational-point codes and natural precision
indices; their outputs are rational vector or matrix codes.
The following conditions are universally quantified numerical
formulas. Their matrices are split into the corresponding
blocks \(\widehat U,\widehat V\).

\paragraph{R0: coherence at rational points.}
For each rational \(q\in\Omega\) and all \(t,u\),
\[
 \norm{\Hhat(q,t)-\Hhat(q,u)}\le2^{-t}+2^{-u},\qquad
 \norm{\Jhat(q,t)-\Jhat(q,u)}\le2^{-t}+2^{-u}.
\]

\paragraph{R1: derivative bounds.}
For every rational \(q\in\Omega\) and every \(t\),
\[
 \norm{\Id-\widehat V(q,t)}\le\tfrac12+2^{-t},
 \qquad \norm{\widehat U(q,t)}\le M+2^{-t}.
\]

\paragraph{R2: base residual.}
For every \(t\),
\[
 \norm{\Hhat((a,b),t)}\le r/4+2^{-t}.
\]
For the exact-base version replace \(r/4+2^{-t}\) by \(2^{-t}\).

\paragraph{R3: rational Taylor certificates.}
For rational \(q,q'\in\Omega\), let \(D=\norm{q'-q}\).
Whenever \(D\le2^{-\mu_0(p)}\), require, for every \(t\),
\begin{equation}\label{eq:rawtaylor}
 \norm{\Hhat(q',t)-\Hhat(q,t)-\Jhat(q,t)(q'-q)}
 \le 2^{-p}D+2^{-t}(2+D).
\end{equation}
The term \(2^{-t}(2+D)\) is the explicit allowance for the two
value errors and the single full-Jacobian error.

\paragraph{R4: rational derivative continuity.}
If \(\norm{q'-q}\le2^{-\nu_0(p)}\), then for all \(t,u\),
\[
 \norm{\Jhat(q',t)-\Jhat(q,u)}
 \le2^{-p}+2^{-t}+2^{-u}.
\]

All norms here are norms of rational data, hence primitive
recursive rational expressions. The domain guards are finite
coordinate comparisons. The input also includes \(r,\alpha>0\),
\(M\ge0\), and \(M\alpha\le r/4\).
If convenient, replace the raw moduli by their bounded prefix
maxima to make them nondecreasing.

R0 defines values and derivative matrices at rational points.
R1--R3 imply the rational segment estimates by the finite
subdivision proof. In particular \(H\) has the explicit
Lipschitz bound \(K=M+3/2\), so it extends to any named point
\(z\in\Omega\). With \(h=\lgup(\max(1,2K))\), a value name is
\[
 H_z(p)=\Hhat(P_\Omega(z(p+h)),p+1).
\]
Its distance to the represented value is at most
\(K2^{-(p+h)}+2^{-(p+1)}\le2^{-p}\).
For the derivative, one possible name is
\[
 J_z(p)=
 \Jhat(P_\Omega(z(\nu_0(p+2)+1)),p+1).
\]
R4 bounds the change to the limiting point by \(2^{-(p+2)}\);
the evaluation error contributes \(2^{-(p+1)}\).
These estimates, expressed pairwise for fine rational
approximants, prove coherence and independence of the input
name.

Safe moduli for named points are
\begin{equation}\label{eq:rawenlarge}
 \mu(p)=\mu_0(p+1)+1,\qquad
 \nu(p)=\nu_0(p+1)+1.
\end{equation}
For example, suppose
\(\norm{z'-z}\le2^{-\mu_0(p+1)-1}\).
Fine clamped rational approximants \(q,q'\) have combined
endpoint error at most another
\(2^{-\mu_0(p+1)-1}\). They therefore satisfy the raw
distance side condition at precision \(p+1\).
Apply \eqref{eq:rawtaylor}, and make every approximation error
smaller than the requested observation error. The resulting
remainder is bounded by \(2^{-(p+1)}\norm{z'-z}\), hence
by \(2^{-p}\norm{z'-z}\).
The same argument proves the derivative-continuity contract
with the enlarged \(\nu\). Explicitly, endpoint-name indices
at least \(\mu_0(p+1)+2\), or respectively
\(\nu_0(p+1)+2\), already ensure the distance guard; the
remaining indices are increased using the displayed
continuity and arithmetic error bounds.

This explains the harmless margin in \eqref{eq:rawenlarge}.
Without a margin, a closed distance guard might fail for
every finite approximation to a pair lying exactly on its
boundary.

\subsection{Removing approximation slack constructively}

\begin{lemma}[Rational squeezing]\label{lem:squeeze}
For rational \(a,b,C\) with \(C\ge0\),
\[
 \bigl(\forall t,\ a\le b+C2^{-t}\bigr)
 \quad\Longrightarrow\quad a\le b
\]
is derivable in \(\ipra\).
\end{lemma}

\begin{proof}
Decide the rational comparison \(a\le b\).
If it fails, \(d=a-b\) is a positive rational.
If \(C=0\), the premise at \(t=0\) is already contradictory.
Otherwise choose \(t=\lgup(2C/d)\); it is a bounded
primitive recursive calculation and gives
\(C2^{-t}\le d/2\).
The premise would then give \(d\le d/2\), contradicting
\(d>0\). This is decidable rational reasoning in intuitionistic
logic. It is not a use of excluded middle for real equality
or of Markov's principle.
\end{proof}

For represented inequalities the goal is itself a universal
collection of rational observations. Fix the requested
observation first, allocate its positive error among the
finitely many quantities, and apply the rational result when
an exact rational comparison is needed. Geometric error
terms are eliminated by the explicit \(\lgup\) shift.
Thus expressions such as allowing a precision to increase
are convenient descriptions of a known finite substitution,
not appeals to an infinitary inference rule.

\subsection{Exactly where induction is used}

There are five kinds of induction in the proof.
\begin{enumerate}
 \item \emph{Finite list arithmetic.} Prove sums, row norms,
 matrix products, and polynomial telescoping identities by
 induction on a finite list length or index.
 \item \emph{Finite polygon sums.} With target precision and
 mesh length fixed, the invariant says that the accumulated
 rational signed error is bounded by the sum of the
 allotted rational local errors.
 \item \emph{Trace preservation.} Each projected rational
 iterate lies in the box, a finite conjunction of rational
 comparisons.
 \item \emph{Inexact recurrence.} The invariant is
 \eqref{eq:scaled-inv}, with the target two traces and their
 precisions fixed.
 \item \emph{Finite matrix loops.} The identities for the
 truncated geometric sum and the telescoping difference of
 powers are componentwise rational equalities.
\end{enumerate}
Every assertion in this list is quantifier-free after coding.
For a finite coordinate or prefix conjunction use a primitive
recursive Boolean fold and the assertion that the fold equals
one. Its extension by one entry is a defining equation.
A universal statement about all precision indices is never
the induction formula.

For a completely literal natural-number version of
Lemma~\ref{lem:recurrence}, fix the full target trace first.
Let \(D_*\) be a positive product of denominators of all
\(d_0,\ldots,d_N,\eta\). This finite product is primitive
recursive. Set \(a_i=D_*d_i\), \(e=D_*\eta\).
These are nonnegative integers, and the invariant is
\[
 2^i a_i+2e\le a_0+2^{i+1}e.
 \tag{Integer invariant}
\]
\(D_*\) is fixed while the prefix index \(i\) varies.
This avoids changing the common denominator inside the
induction. The compiled Lean core proves this integer
recurrence and a finite Boolean-checker form of it.

The contraction inequality needed at a particular step is
obtained from the raw Taylor certificates by the polygon
argument and rational squeezing. The resulting distance
between the two actual program outputs is rational.
Consequently the recurrence premise itself is a
quantifier-free rational inequality, even though its
derivation consulted universally quantified input
certificates.

\subsection{Finite matrix estimates at an observation precision}

The same discipline applies to the Neumann argument.
Fix the finite power index \(N\), and approximate the
represented coefficient matrix at a rational precision.
Induct on products of those rational matrices. A norm
bound \(\norm{\widetilde E}\le1/2+\delta\) gives
\[
 \norm{\widetilde E^N}\le(1/2+\delta)^N
\]
by rational induction. The finite identity
\[
 u^N-v^N=(u-v)\sum_{i=0}^{N-1}u^i v^{N-1-i}
\]
bounds \((1/2+\delta)^N-2^{-N}\) by
\(N\delta(3/4)^{N-1}\) when \(N\ge1\) and
\(\delta\le1/4\). Choosing \(\delta\) from \(N\) and the
requested observation error removes this difference.
For \(N=0\) the identity is exact.
Finite sums and products of matrix entries have equally
explicit error propagation bounds. Lemma~\ref{lem:perturb}
gives the sharper uniform constants used for the actual
inverse and derivative programs.

Thus the finite algebra is first proved for rational data
and transferred to named coefficients at fixed finite
indices. There is no hidden induction over a statement
quantifying over all observations of a real-matrix power.

\subsection{The precise theorem schema}

Fix primitive recursive definitions for the input
evaluators and moduli. Adjoin the derived programs
\(\Ghat,\widehat D\) by their primitive recursive defining
equations. Let \(\operatorname{Hyp}\) denote R0--R4 and the
numerical box conditions. The proof above constructs an
iPRA proof of
\[
 \operatorname{Hyp}\ \longrightarrow\
 \bigl(\operatorname{Coherent}(\Ghat)
 \wedge\operatorname{Zero}(\Ghat)
 \wedge\operatorname{Unique}(\Ghat)
 \wedge\operatorname{Lip}_{2M}(\Ghat)
 \wedge\operatorname{Taylor}_{\mu_g}(\Ghat,\widehat D)
 \wedge\operatorname{UC}_{\nu_g}(\widehat D)\bigr).
 \tag{IFT schema}
\]
The displayed right side is notation for the numerical
observation formulas described above, universally closed
over rational inputs, precision indices, and any numerical
parameters of the fixed presentation. For a fixed additional
input-point name, substitute the explicit name-composition
program. This is a metatheorem about a family of first-order
proofs, not one object theorem quantifying over a type of
all primitive recursive functions.

If proofs of R0--R4 are supplied in iPRA, modus ponens
discharges the hypotheses. The polynomial examples below
have direct finite algebraic proofs of these hypotheses.
Numerical testing of many input points is not a substitute
for these universal proofs.

The handwritten argument specifies each infinite-looking
step by finite observation and each induction by a
quantifier-free invariant. It is a proof at the requested
human-readable level. Constructing and checking the entire
corresponding object-language proof term remains a separate
formalization task, specified in Section~\ref{sec:lean}.
